\section{Execution and Repair}
Now, we move from correctness-checking to code understanding as a whole. First, we evaluate the execution prediction and repair capabilities of language models. Then, we compare them to a model's code generation capability, answering the question: \textit{are models better at understanding their generated programs for problems they can confidently solve}?
\txo{Suggestion: to distinguish this from the previous section a bit more, maybe phrase it like \textit{are models better at predicting the semantics of programs for tasks that they can reliably solve?} (point is: the last section also tests their ability to "understand", so hone in on what makes the "understanding" in this section different than the previous one}

\textbf{Can models execute code they generate?} Next, we measure the ability of models to execute model-generated programs. We report accuracies with CoT for GPT-3.5 and GPT-4 and without CoT for other models, as we find they do not benefit from CoT on this task. In Fig. \ref{fig:execution-accuracy} (left), we show the execution accuracies of six open models on LeetCode problems (others in Fig. \ref{fig:execution-results-full}), observing accuracies between 0-40\%.

\textbf{How do models behave when asked to execute incorrect code?}
\txo{I personally feel like this section and the previous one should be merged (and cut a bit).}
Next, we partition the samples in the dataset based on whether or not the program is correct for the given execution input, shown in Fig. \ref{fig:execution-accuracy} (right) for HumanEval programs (others in Fig. \ref{fig:execution-breakdown-full}). The accuracies for correct and incorrect samples are shown in green and red, respectively, showing that models are much better at predicting the output of correct programs than incorrect ones. For the incorrect programs, we find that models other than GPT-4 frequently predicted the output that would have resulted had the program been correct (light blue) just as often as they predicted the correct output. Since models only see the programs and not the problem statements, this suggests that models may frequently hallucinate an incorrect program as a correct one, consistent with our previous finding that models cannot reliably assess the correctness of programs.

\begin{figure}[ht!]
    \centering
    \begin{subfigure}[b]{0.49\columnwidth}
        \centering
        \includegraphics[width=\textwidth]{figures/accuracy_barplots/lcb_execution.pdf}
    \end{subfigure}
    \hfill
    \begin{subfigure}[b]{0.49\columnwidth}
        \centering
        \includegraphics[width=\textwidth]{figures/execution_breakdown/humaneval_codellama-34b_sample.pdf}
    \end{subfigure}
    \caption{Results on execution prediction task. Left: accuracies of 6 models on LeetCode programs. Right: accuracies on HumanEval, showing that models do better executing correct programs than incorrect ones and often execute as if the program were correct.}
    \label{fig:execution-accuracy}
\end{figure}

\textbf{(Q3) To what extent can the LMs identify and repair programs they generate, without any external feedback?}


\txo{TODO}




\textbf{(Q4) How does problem difficulty influence a model's ability to perform correctness checking, execution, and repair on programs of that problem?} One might expect that a model might have a better understanding of programs that corresponding to ``easier'' problems, i.e. problems it can easily generate correct solutions to. To test this, for a given model, we bucket problem difficulties into easy, medium, and hard by checking how likely the model is to generate a correct solution to that problem, i.e. its pass@1 on that problem. We then calculate the average accuracy for each of the tasks for each problem difficulty bucket. In Fig. \ref{fig:difficulty-sample}, we show a few of these results on HumanEval, deferring with the full set to Sec. \ref{appendix:difficulty-correctness-full}, \ref{appendix:difficulty-execution-full}, and \ref{appendix:difficulty-repair-full}. We find that 1) correctness checking accuracies are relatively uncorrelated with problem difficulty, 2) execution ability is better for easier problems, and 3) repair ability is ??. These trends are robust and generally hold across HumanEval, LeetCode, and ODEX and many models. 

\begin{figure*}[ht!]
    \centering
    \begin{subfigure}[b]{0.31\textwidth}
        \centering
        \includegraphics[width=\textwidth]{figures/difficulty_barplots/humaneval_correctness_sample.pdf}
        \caption{Correctness Checking}
    \end{subfigure}
    \hfill
    \begin{subfigure}[b]{0.31\textwidth}
        \centering
        \includegraphics[width=\textwidth]{figures/difficulty_barplots/humaneval_execution_sample.pdf}
        \caption{Execution}
    \end{subfigure}
    \hfill
    \begin{subfigure}[b]{0.31\textwidth}
        \centering        \includegraphics[width=\textwidth]{figures/misc/mink.jpg}
        \caption{Repair}
    \end{subfigure}
    \caption{Accuracies of correctness checking, execution, and repair on HumanEval. A problem's difficulty is determined by the pass@1 of the data-generating model. In general, correctness checking accuracies are relatively uncorrelated with problem difficulty, execution of easier problems is easier for models, and...}
    \label{fig:difficulty-sample}
\end{figure*}

% \begin{figure}
%     \centering
%     \includegraphics[width=\columnwidth]{figures/execution_breakdown/humaneval_codellama-34b_sample.pdf}
%     \caption{Models are better at executing correct programs than incorrect ones. In addition, they often execute incorrect programs as if they were correct.}
%     \label{fig:execution-breakdown-sample}
% \end{figure}
% \textbf{Code Verification}: 
